\chapter{Seven Strategic Challenges for Contemporary Europe}

\begingroup

\begin{abstract}
	This article develops a concise typology of seven strategic challenges facing the European Union and the wider European political order. It treats these challenges not as identical threats, but as pressures operating through different mechanisms: military coercion, economic dependency, alliance volatility, domestic nationalism, illiberal currents within the left, religiously framed political radicalism, and the fragility of democratic governance itself. The argument is that Europe's survival as a liberal democratic project depends on its capacity to distinguish among these pressures while strengthening its own institutional, military, economic, and civic resilience.
\end{abstract}

\section{Introduction}

The European Union and contemporary Europe more broadly face a strategic environment marked by external coercion, internal fragmentation, and uncertainty about the reliability of inherited alliances. These pressures are not reducible to a single adversary or ideological family. They arise from state actors, transnational movements, domestic parties, social cleavages, and weaknesses internal to democratic government. A political science account therefore needs to separate the mechanisms through which each pressure operates. The following study identifies seven challenges: Russia, China, the United States under illiberal leadership, nationalist movements inside Europe, destabilizing tendencies within parts of the left, political Islamism, and the demanding maintenance of democracy itself.

\section{Russia and coercive revisionism}

The first challenge is the Russian Federation's revisionist and imperial posture toward its neighbors. Russia's war against Ukraine is the clearest expression of a wider pattern: the use of historical grievance, security pretexts, disinformation, energy leverage, cyber operations, and conventional force to contest the sovereignty of nearby states. Although the institutions and ideology of Tsarist Russia, the Soviet Union, and the Russian Federation differ, European security policy must recognize a recurring strategic habit: Moscow's tendency to treat the autonomy of neighboring peoples as conditional. This makes Russia the most immediate military threat to Europe. It requires deterrence, sustained support for threatened states, resilience against hybrid operations, and a European defense capacity that does not rely entirely on external guarantees.

\section{China and strategic dependency}

China presents a different challenge. It is less immediate in military terms for Europe, but potentially more consequential in economic and infrastructural terms. The issue is not the presence of Chinese communities in Europe, which should be treated within ordinary democratic standards of equal citizenship and integration policy. The strategic concern lies elsewhere: state-linked investment, control of logistics and communications infrastructure, technology dependency, market access asymmetries, and the political leverage that can follow from economic concentration. Ports, telecommunications networks, industrial supply chains, and critical raw materials are all relevant sites of vulnerability. Because China combines major economic capacity with authoritarian governance, Europe must protect openness without allowing openness to become dependency.

\section{Alliance volatility and the United States}

The United States remains a central actor in European security, but recent American politics has exposed the danger of excessive dependence on a single ally. When an allied democracy moves toward oligarchic politics, institutional erosion, or unilateral foreign policy, its partners face strategic uncertainty even without becoming formal enemies. For Europe, the problem is not anti-Americanism. It is the recognition that the transatlantic alliance can no longer substitute for European political agency. Sudden shifts in trade policy, military commitments, diplomatic reliability, and attitudes toward authoritarian powers can impose severe costs on European states. A mature European strategy should preserve cooperation with the United States while reducing vulnerability to American domestic turbulence.

\section{Nationalist movements within Europe}

A fourth challenge comes from nationalist and populist movements inside the European Union. The Union is built on the partial pooling of sovereignty among member states. Parties that define sovereignty in strictly national and zero-sum terms therefore tend to weaken common action in foreign policy, fiscal coordination, migration governance, energy security, and defense. Not every form of national attachment is anti-European, and democratic politics must leave room for legitimate disputes over integration. The strategic problem arises when nationalism becomes a systematic project to paralyze collective institutions while benefiting from them. Europe's international influence depends on scale, coordination, and credibility; nationalist obstruction reduces all three.

\section{Illiberal currents within the Left}

A fifth challenge is symmetrical in part, though not identical: currents within the left that subordinate liberal democratic institutions to abstract internationalism, anti-Western reflexes, or an uncritical politics of open borders. Left traditions have also defended rights, labor protections, anti-authoritarianism, and social citizenship, so the problem cannot be attributed to the left as such. The concern is a narrower tendency: when parts of the left treat liberal democracy as merely a mask for domination, they may become inattentive to the institutional conditions that make dissent, minority rights, and social reform possible. Migration policy is one field where this tension appears sharply. Europe needs humane asylum and integration policies, but large-scale migration without civic integration, administrative capacity, and democratic consensus can produce social fragmentation and strengthen illiberal backlash.

\section{Political islamism and liberal order}

The sixth challenge concerns political Islamism, understood not as Islam as a faith or Muslims as citizens, but as movements that seek to subordinate political authority, law, gender equality, and individual rights to a broad religious-political project. Europe must distinguish carefully between religious liberty and ideological movements hostile to pluralist constitutionalism. The historical record of religious expansionism is not unique to Islam; Christianity also produced imperial and coercive forms that Europe has had to overcome through secularization, rights, and constitutional limits on clerical power. The relevant question today is whether European institutions can protect Muslim citizens, defend freedom of conscience, and resist organizations or networks that reject democratic equality. Confusing these categories weakens both security and liberal legitimacy.

\section{Democracy as a permanent challenge}

The final challenge is democracy itself. Democracy is not only an institutional arrangement but an ongoing practice that requires maintenance. Separation of powers, independent courts, free journalism, public education, civic trust, party competition, and the formation of informed public opinion all need constant protection. As political systems grow larger and more complex, citizens find it harder to understand where power lies and how decisions are made. This can create alienation, conspiratorial thinking, and openness to authoritarian simplification. The European Union faces this problem acutely because it combines national democracies with supranational institutions. Its legitimacy depends on making shared rule intelligible, accountable, and capable of action.

\section{Conclusion}

Europe's strategic predicament is multidimensional. Russia threatens through military coercion and hybrid warfare. China threatens through dependency and authoritarian economic leverage. The United States, even as an ally, can become a source of uncertainty when its domestic politics undermine reliability. Nationalist and illiberal movements weaken Europe's internal capacity for collective action. Political Islamism tests the boundary between religious freedom and constitutional order. Democracy itself remains fragile because it depends on habits, institutions, and knowledge that cannot be secured once and for all. The central task is not simply to identify enemies. It is to build a Europe capable of defending liberal democracy without abandoning the principles that make it worth defending.

\endgroup
